\chapter{Collision Detection for Non-Convex Polygons}

Evaluating the energy function $E(\conf; L)$ at every Markov chain step
requires determining, for each pair of polygons, whether they overlap
and by how much.
This computation is the dominant cost kernel of the entire optimization:
for $N$ polygons, there are $\binom{N}{2} = O(N^2)$ candidate pairs,
and each narrowphase test involves nontrivial geometric work.
The collision system occupies the same structural role in this project
as force evaluation in molecular dynamics or matrix multiplication in
linear algebra—its throughput sets the ceiling on how many SA steps
can be executed per second.

This chapter develops the mathematical foundations of collision detection
for non-convex polygons, presents the three narrowphase implementations
benchmarked in this study, and analyzes their cost structure.

%─────────────────────────────────────────────────────────────────
\section{Non-Convex Polygons and Convex Decomposition}
%─────────────────────────────────────────────────────────────────

The polygon used in this study is a planar non-convex (concave) polygon
$\mathcal{P} \subset \mathbb{R}^2$ with $V$ vertices.
A polygon is convex if for every pair of interior points $p, q$,
the segment $[p, q]$ lies entirely within $\mathcal{P}$.
Non-convex polygons violate this condition; equivalently, their
boundary is not everywhere locally on the same side of every edge
normal.

The Separating Axis Theorem, which provides the most efficient exact
intersection test, applies only to convex shapes.
To handle non-convex polygons, we decompose $\mathcal{P}$ into a
collection of convex \emph{parts}:
\[
  \mathcal{P} = \bigcup_{j=1}^{M} C_j, \qquad
  C_j \text{ convex},
\]
where the parts cover $\mathcal{P}$ exactly (with at most boundary overlap).
Two non-convex polygon instances $\mathcal{P}_i$ and $\mathcal{P}_k$
(at poses $(x_i, y_i, \theta_i)$ and $(x_k, y_k, \theta_k)$) intersect
if and only if at least one part of $\mathcal{P}_i$ intersects at least
one part of $\mathcal{P}_k$:
\[
  \mathcal{P}_i \cap \mathcal{P}_k \neq \emptyset
  \;\iff\;
  \exists\, j \in \{1,\dots,M\},\; l \in \{1,\dots,M\}
  \;\text{ s.t. }\;
  C_j^{(i)} \cap C_l^{(k)} \neq \emptyset,
\]
where $C_j^{(i)}$ denotes part $j$ of polygon $i$ transformed to
world space. This reduces non-convex intersection to a collection of
convex--convex tests.

\paragraph{Decomposition methods.}
Two decompositions are used in this study:
\begin{itemize}
  \item \textbf{Triangulation (Method A).}
        $\mathcal{P}$ is decomposed into $M = V - 2$ triangles via
        fan triangulation from a reference vertex. Each part $C_j$ is
        a triangle, giving the smallest possible number of vertices per
        part (3) at the cost of the largest number of parts.
  \item \textbf{Convex decomposition (Methods B and C).}
        $\mathcal{P}$ is decomposed into a smaller number of convex
        polygons, each with more vertices. Fewer parts reduce the
        number of part--part pairs; larger vertex counts increase the
        per-pair axis budget.
\end{itemize}

%─────────────────────────────────────────────────────────────────
\section{Coordinate Transforms}
%─────────────────────────────────────────────────────────────────

Each polygon instance $i$ is defined by a pose
$(x_i, y_i, \theta_i) \in \mathbb{R}^2 \times [0, 2\pi)$.
Local-frame vertices $v \in \mathbb{R}^2$ are mapped to world frame by
the rigid-body transform
\[
  T_i(v) = R(\theta_i)\, v + \mathbf{t}_i,
  \qquad
  R(\theta) =
  \begin{pmatrix} \cos\theta & -\sin\theta \\
                  \sin\theta & \phantom{-}\cos\theta \end{pmatrix},
  \quad
  \mathbf{t}_i = (x_i, y_i)^\top.
\]
The world-frame vertices of part $C_j^{(i)}$ are
$\{T_i(v) : v \in C_j\}$.
Edge normals transform under rotation only: if $\mathbf{n}$ is a
unit outward normal in local frame, its world-frame counterpart is
$R(\theta_i)\,\mathbf{n}$.

A naive implementation recomputes $R(\theta_i)$ and all transformed
vertices inside the innermost pair loop, evaluating $\sin$ and $\cos$
once per pair--part test. The cached method (Section~\ref{sec:method-c})
amortizes this cost over all tests involving polygon $i$.

%─────────────────────────────────────────────────────────────────
\section{The Separating Axis Theorem}
\label{sec:sat}
%─────────────────────────────────────────────────────────────────

\subsection{Statement and Proof Sketch}

\begin{theorem}[Separating Axis Theorem]
Let $P, Q \subset \mathbb{R}^2$ be compact convex sets.
Then $P \cap Q = \emptyset$ if and only if there exists a unit vector
$\mathbf{n} \in \mathbb{R}^2$ such that
\[
  \max_{p \in P}\, \langle p, \mathbf{n} \rangle
  \;<\;
  \min_{q \in Q}\, \langle q, \mathbf{n} \rangle.
\]
Such a vector is called a \emph{separating axis}.
\end{theorem}

The ``if'' direction is immediate: if a separating axis exists then no
point belongs to both projected intervals, so $P \cap Q = \emptyset$.
The ``only if'' direction follows from the separating hyperplane theorem
for convex sets: if $P \cap Q = \emptyset$, by convexity and compactness
there exists a hyperplane (in $\mathbb{R}^2$, a line) strictly separating
them, whose normal is the required $\mathbf{n}$.

\subsection{Projection and Interval Test}

For a convex polygon $P$ with vertices $\{v_1, \dots, v_k\}$, the
projection onto axis $\mathbf{n}$ yields the interval
\[
  I_P(\mathbf{n}) = \Bigl[
    \min_{1 \le j \le k} \langle v_j, \mathbf{n} \rangle,\;
    \max_{1 \le j \le k} \langle v_j, \mathbf{n} \rangle
  \Bigr].
\]
Two intervals $I_P = [a, b]$ and $I_Q = [c, d]$ are \emph{disjoint}
if $b < c$ or $d < a$. Disjointness on any axis certifies non-intersection.

\subsection{Candidate Axes for Polygons}

For two convex polygons with $k_1$ and $k_2$ vertices respectively,
it suffices to test only $k_1 + k_2$ candidate axes: the outward edge
normals of each polygon. This follows from the observation that the
extremal separating direction between two convex polygons must be
perpendicular to a face of one of them.

\begin{corollary}
Two convex polygons $P$ (with $k_1$ edges) and $Q$ (with $k_2$ edges)
are disjoint if and only if a separating axis exists among the $k_1 + k_2$
outward edge normals of $P$ and $Q$.
\end{corollary}

The algorithm tests axes in sequence and returns \textsc{disjoint} as
soon as any separating axis is found (\emph{early-out}).
If all $k_1 + k_2$ axes yield overlapping intervals, the algorithm
returns \textsc{intersecting}.

\subsection{Penetration Depth}

When the polygons do intersect, the minimum translational distance (MTD)
required to separate them is the minimum interval overlap over all axes:
\[
  d_{\mathrm{pen}} = \min_{\mathbf{n} \in \mathcal{A}}
  \max\!\bigl(0,\;
    \min(b_{P}, b_{Q}) - \max(a_{P}, a_{Q})
  \bigr),
\]
where $I_P(\mathbf{n}) = [a_P, b_P]$, $I_Q(\mathbf{n}) = [a_Q, b_Q]$,
and $\mathcal{A}$ is the set of candidate axes.
The MTD provides a natural penalty signal: deeper interpenetrations
contribute more to the energy function.

%─────────────────────────────────────────────────────────────────
\section{Broadphase Filtering}
\label{sec:broadphase}
%─────────────────────────────────────────────────────────────────

For $N$ polygons, there are $\binom{N}{2}$ instance pairs and
$\binom{NM}{2}$ part pairs in the worst case.
Narrowphase SAT is applied only to pairs that pass a cheap
\emph{broadphase} filter, which conservatively rejects pairs
that cannot possibly intersect.

\subsection{Axis-Aligned Bounding Box Test}

The axis-aligned bounding box (AABB) of a transformed polygon instance
$\mathcal{P}_i$ in world frame is
\[
  \mathrm{AABB}(\mathcal{P}_i)
  = \Bigl[\min_j (T_i(v_j))_x,\; \max_j (T_i(v_j))_x\Bigr]
  \times
  \Bigl[\min_j (T_i(v_j))_y,\; \max_j (T_i(v_j))_y\Bigr].
\]
Two polygons can intersect only if their AABBs overlap, since the AABB
is a convex superset of the polygon. AABB overlap is tested in $O(1)$
via four scalar comparisons and rejects non-proximate pairs with no
trigonometric cost.

Per-part AABBs are computed analogously for each convex part $C_j^{(i)}$,
providing a second level of filtering before the narrowphase is invoked.

\subsection{Spatial Hash Grid}

For $N > 10$, the cost of testing all $\binom{N}{2}$ AABB pairs
becomes significant. A spatial hash grid partitions the domain
$[0, L]^2$ into cells of width $h$ and assigns each polygon to the
cells overlapping its AABB. Only polygons sharing at least one cell
are candidate pairs.

The grid is stored as a compact hash table: cell keys are hashed to
buckets, and each bucket stores a linked list of polygon indices.
For a uniform random configuration at packing fraction $\phi$, the
expected number of neighbors per polygon is $O(1)$ for a cell size
$h$ proportional to the polygon diameter, giving $O(N)$ total candidate
pairs and reducing the broadphase from $O(N^2)$ to $O(N)$ in expectation.

The grid is rebuilt every time a configuration is modified---captured
in the \texttt{rebuild\_derived} call---so its cost is amortized over
all pair tests in the subsequent energy evaluation.

%─────────────────────────────────────────────────────────────────
\section{Narrowphase Method A: Triangle Decomposition}
\label{sec:method-a}
%─────────────────────────────────────────────────────────────────

Method A decomposes each polygon into $M = V - 2$ triangles via fan
triangulation from vertex $v_1$:
\[
  \mathcal{P} = \bigcup_{j=2}^{V-1} \triangle(v_1, v_j, v_{j+1}).
\]
Intersection between two polygon instances is tested by checking all
$M^2 = (V-2)^2$ triangle--triangle pairs.

\paragraph{Triangle--triangle test.}
Two triangles $\triangle ABC$ and $\triangle DEF$ intersect if and only
if any edge of one separates them under SAT with 6 candidate axes
(3 edge normals per triangle), or one triangle contains a vertex of
the other. This is a fully degenerate instance of the general convex
polygon SAT with $k_1 = k_2 = 3$.

\paragraph{Cost.}
Each triangle--triangle test requires at most $6$ axis projections,
each of cost $O(3)$ dot products, giving $O(18)$ floating-point
operations per pair in the worst case with no early-out.
The total narrowphase cost for one polygon--polygon test is
$O(M^2) = O((V-2)^2)$ in the worst case, with early-out reducing
average cost substantially in non-intersecting configurations.

%─────────────────────────────────────────────────────────────────
\section{Narrowphase Method B: Baseline Convex SAT}
\label{sec:method-b}
%─────────────────────────────────────────────────────────────────

Method B applies SAT directly to each pair of convex parts
$C_j^{(i)}$ and $C_l^{(k)}$ in world coordinates.
For each part pair, the algorithm:
\begin{enumerate}
  \item Recomputes the world-frame vertices of $C_j^{(i)}$ by applying
        $T_i$ to all local vertices.
  \item Recomputes the world-frame edge normals of $C_j^{(i)}$ by
        applying $R(\theta_i)$ to all local normals.
  \item Repeats steps 1--2 for $C_l^{(k)}$.
  \item Tests all $k_j + k_l$ candidate axes via projection and
        interval overlap.
\end{enumerate}
Steps 1--3 are performed fresh for every part pair, even when polygon
$i$ appears in multiple tests during the same energy evaluation.
If polygon $i$ has $D$ neighbors, its transform $T_i$ is applied
$D \cdot M$ times: once per part per neighbor.

\paragraph{Cost.}
For a polygon with $M$ parts each of $k$ vertices, evaluating
$E_{\mathrm{overlap}}$ for one polygon against all $D$ neighbors costs
\[
  O(D \cdot M^2 \cdot k)
\]
floating-point operations, where the $M^2$ factor counts part--part
pairs and the $k$ factor counts projections per axis.
The trigonometric cost ($\sin$, $\cos$) is $O(D \cdot M)$.

%─────────────────────────────────────────────────────────────────
\section{Narrowphase Method C: Cached Convex SAT}
\label{sec:method-c}
%─────────────────────────────────────────────────────────────────

Method C eliminates the redundant per-pair transform recomputations
by caching world-frame data for all $N$ polygons in a per-thread
\texttt{Workspace} before the pair loop begins.

\subsection{Pre-Pass: Cache Population}

Before any pair test, a single pass over all $N$ polygons computes and
stores in the \texttt{Workspace}:
\begin{align}
  \text{World vertices:}\quad &
  v^{(i)}_{\mathrm{world}} = T_i(v^{(i)}_{\mathrm{local}})
  \quad \forall i, \forall v, \label{eq:cache-verts} \\
  \text{World normals:}\quad &
  \mathbf{n}^{(i)}_{\mathrm{world}} = R(\theta_i)\,\mathbf{n}^{(i)}_{\mathrm{local}}
  \quad \forall i, \forall \text{ edge}, \label{eq:cache-normals} \\
  \text{AABBs:}\quad &
  \mathrm{AABB}(C_j^{(i)})
  \quad \forall i, \forall j. \label{eq:cache-aabb}
\end{align}
The rotation $R(\theta_i)$ is evaluated once per polygon per energy
evaluation, with $\sin(\theta_i)$ and $\cos(\theta_i)$ stored as scalars.
All subsequent axis projections read from the cached arrays.
Cache population costs $O(NMk)$ uniformly, amortized as $O(k/D)$ per
pair test for a polygon with $D$ neighbors.

\subsection{Hint-Axis Pre-Filter}

Before executing the full $k_j + k_l$-axis SAT loop, Method C tests a
single \emph{hint axis}: the unit vector connecting the centroids of
the two parts:
\[
  \mathbf{n}_{\mathrm{hint}}
  = \frac{\bar{C}_j^{(i)} - \bar{C}_l^{(k)}}
         {\|\bar{C}_j^{(i)} - \bar{C}_l^{(k)}\|},
\]
where $\bar{C}_j^{(i)}$ denotes the centroid of part $j$ of instance $i$.
If the projections of $C_j^{(i)}$ and $C_l^{(k)}$ onto
$\mathbf{n}_{\mathrm{hint}}$ are disjoint, the parts are immediately
classified as non-intersecting at the cost of $2k$ dot products---
substantially cheaper than the full SAT loop.

The hint axis is effective because separated polygons in a packing
are most commonly separated along the direction between their centers.
It provides a ``free'' early-out for the majority of non-intersecting
pairs that survive AABB filtering.

\subsection{Full SAT Fallback}

If the hint axis does not separate the parts, the full SAT loop over
all $k_j + k_l$ cached world normals is executed with early-out.
Projections read directly from the pre-computed vertex arrays
(\ref{eq:cache-verts})--(\ref{eq:cache-normals}), avoiding all
trigonometric recomputation.

\subsection{Cost Reduction}

Let $D$ be the average number of polygon neighbors, $M$ the parts per
polygon, and $k$ the vertices per part.
Table~\ref{tab:cost} compares the operation counts per energy evaluation.

\begin{table}[h]
\centering
\caption{Approximate floating-point operation counts per energy evaluation.}
\label{tab:cost}
\begin{tabular}{lll}
\toprule
Step & Method B & Method C \\
\midrule
$\sin$/$\cos$ evaluations & $O(DNM)$ & $O(N)$ \\
Vertex transforms & $O(DNMk)$ & $O(NMk)$ (pre-pass) \\
Axis projections (per pair) & $O(M^2 k)$ & $O(k)$ hint + $O(M^2 k)$ fallback \\
\bottomrule
\end{tabular}
\end{table}

For $D \gg 1$ (dense packing) the trigonometric saving is the dominant
effect: Method C evaluates $\sin$/$\cos$ exactly $N$ times per energy
call regardless of density, while Method B evaluates them $DNM$ times.

%─────────────────────────────────────────────────────────────────
\section{Connection to the Energy Function}
\label{sec:energy-collision}
%─────────────────────────────────────────────────────────────────

The collision system feeds directly into the energy function introduced
in Chapter~1:
\[
  E(\conf; L) = E_{\mathrm{overlap}}(\conf) + E_{\mathrm{boundary}}(\conf; L).
\]

\paragraph{Overlap penalty.}
For each pair of polygon instances $(i, k)$ found to intersect by the
narrowphase, the overlap penalty is the sum of penetration depths over
all intersecting part pairs:
\[
  E_{\mathrm{overlap}}(\conf)
  = \sum_{i < k} \sum_{\substack{j,\, l \\ C_j^{(i)} \cap C_l^{(k)} \neq \emptyset}}
    d_{\mathrm{pen}}\!\left(C_j^{(i)},\, C_l^{(k)}\right).
\]
This penalty is zero if and only if no two polygons overlap, i.e., the
configuration is collision-free.

\paragraph{Boundary penalty.}
For each polygon $i$, the boundary penalty accumulates for every
vertex lying outside the container $[0, L]^2$:
\[
  E_{\mathrm{boundary}}(\conf; L)
  = \sum_{i=1}^N \sum_{v \in \mathcal{P}_i}
    \bigl\| \max(\mathbf{0},\, -T_i(v)) \bigr\|
    + \bigl\| \max(\mathbf{0},\, T_i(v) - L\,\mathbf{1}) \bigr\|,
\]
where the $\max$ is taken componentwise.
A configuration is feasible (Section~\ref{sec:bisection-integration})
when both penalties are below the threshold
$\varepsilon_{\mathrm{feas}}$.

%─────────────────────────────────────────────────────────────────
\section{Benchmark Design}
\label{sec:bench}
%─────────────────────────────────────────────────────────────────

\subsection{Objective}

The microbenchmark isolates the collision system from the SA loop
to measure throughput in controlled conditions. Its purpose is twofold:
(i) to verify that Method C achieves a concrete speedup over Methods A
and B before committing to the 4-hour hero runs, and (ii) to quantify
the rejection efficiency of the broadphase and hint-axis filter, which
determines what fraction of SA steps incur full SAT cost.

\subsection{Protocol}

A fixed set of $N$ polygon instances is placed at random non-overlapping
positions. The full $\binom{N}{2}$ pair evaluation loop is then executed
$R_{\mathrm{bench}}$ times, and wall-clock time is recorded. The same
configuration is used for all three methods, so differences in timing
reflect only the collision algorithm.

The benchmark is invoked as:
\begin{verbatim}
  ./bench <N> <R_bench> <out.csv>
\end{verbatim}

\subsection{Primary Metrics}

The following quantities are recorded for each method:

\begin{table}[h]
\centering
\caption{Benchmark metrics collected per method.}
\label{tab:bench-metrics}
\begin{tabular}{ll}
\toprule
Metric & Definition \\
\midrule
\textbf{Pairs total} & $\binom{N}{2}$ instance pairs evaluated \\
\textbf{Broadphase pass} & Pairs surviving instance AABB filter \\
\textbf{Narrow calls} & Narrowphase invocations \\
\textbf{Collisions} & Intersecting pairs found \\
\textbf{Wall time (s)} & Total elapsed time \\
\textbf{Mpairs/s} & Million instance pairs per second \\
\bottomrule
\end{tabular}
\end{table}

For Method C, additional counters are recorded to expose the filtering
cascade:

\begin{table}[h]
\centering
\caption{Method C internal diagnostic metrics.}
\label{tab:bench-cached}
\begin{tabular}{ll}
\toprule
Metric & Definition \\
\midrule
\textbf{Part pairs} & Total part--part pairs entering narrowphase \\
\textbf{Part AABB pass} & Part pairs surviving per-part AABB filter \\
\textbf{Hint tests} & Hint-axis projections executed \\
\textbf{Hint rejects} & Pairs rejected by hint axis \\
\textbf{SAT full calls} & Full SAT loop invocations \\
\textbf{SAT hits} & SAT calls returning intersection \\
\textbf{SAT axes tested} & Total axes evaluated across SAT calls \\
\textbf{Avg axes / SAT call} & SAT axes tested $\div$ SAT full calls \\
\bottomrule
\end{tabular}
\end{table}

The ratio \textbf{Hint rejects / Part AABB pass} measures the
effectiveness of the hint-axis pre-filter: a high ratio indicates
that most surviving AABB pairs are resolved cheaply without
entering the full SAT loop.
The \textbf{Avg axes / SAT call} metric quantifies early-out
effectiveness within the SAT loop itself.

\subsection{Scaling Study}

The benchmark is run at $N \in \{10, 20, 50, 100, 200\}$ to characterize
how throughput scales with system size. For each $N$, three placement
regimes are evaluated:
\begin{enumerate}
  \item \textbf{Sparse:} polygons placed at random with low overlap
        probability. Tests broadphase efficiency.
  \item \textbf{Dense:} polygons packed at high density near feasibility.
        Tests narrowphase cost in the regime the SA spends most time.
  \item \textbf{Jammed:} polygons at a jammed configuration with maximal
        neighbor count. Tests worst-case narrowphase cost.
\end{enumerate}

\subsection{Results and Interpretation}
\label{sec:bench-results}

\begin{center}
\begin{tabular}{lrrrrr}
\toprule
Method & Pairs total & Broad pass & Narrow calls & Collisions & Time (s) \\
\midrule
Triangle (A) & \texttt{\_} & \texttt{\_} & \texttt{\_} & \texttt{\_} & \texttt{\_} \\
SAT baseline (B) & \texttt{\_} & \texttt{\_} & \texttt{\_} & \texttt{\_} & \texttt{\_} \\
SAT cached (C) & \texttt{\_} & \texttt{\_} & \texttt{\_} & \texttt{\_} & \texttt{\_} \\
\bottomrule
\end{tabular}
\end{center}

\begin{center}
\begin{tabular}{lr}
\toprule
Method C diagnostic & Value \\
\midrule
Part pairs & \texttt{\_} \\
Part AABB pass & \texttt{\_} \\
Hint tests & \texttt{\_} \\
Hint rejects & \texttt{\_} \\
SAT full calls & \texttt{\_} \\
SAT hits & \texttt{\_} \\
SAT axes tested & \texttt{\_} \\
Avg axes / SAT call & \texttt{\_} \\
\bottomrule
\end{tabular}
\end{center}

The expected outcome, borne out by the cascade structure of Method C,
is as follows. Most polygon pairs are rejected by the instance-level
AABB filter, so the broadphase pass rate should be small ($< 20\%$
for sparse configurations). Of the pairs that survive, the per-part
AABB filter and hint-axis pre-filter together reject the majority before
any full SAT loop is invoked. The full SAT loop is called only for pairs
that are either intersecting or geometrically close to intersection,
and among those calls, the early-out terminates the axis loop after
on average far fewer than $k_j + k_l$ tests. This cascade is the
mechanism by which Method C achieves substantially higher throughput
than Methods A and B at equal correctness.

%─────────────────────────────────────────────────────────────────
\section{Integration into the SA Loop}
%─────────────────────────────────────────────────────────────────

In the context of the optimization study, the collision system is
invoked at every accepted SA proposal that changes any polygon's pose.
Per-thread \texttt{Workspace} instances are allocated once at
thread initialization and reused across all energy evaluations,
avoiding heap allocation in the inner loop.

After a pose update, \texttt{rebuild\_derived} recomputes the world
vertices, rotated axes, and AABBs for the moved polygon only, and
updates its spatial hash cell memberships. A full grid rebuild is
triggered only when global configuration changes require it
(warm-start, clone, or swap operations).
This incremental update strategy reduces the pre-pass cost of
Method C from $O(NMk)$ to $O(Mk)$ per SA step in the common case
where only one polygon moves at a time.

The collision throughput therefore sets the wall-clock budget for
SA step rate: if Method C achieves $T_{\mathrm{narrow}}$ microseconds
per narrowphase call and the average polygon has $D$ neighbors at
$M$ parts each, the expected energy evaluation time per SA step is
approximately $D \cdot M^2 \cdot T_{\mathrm{narrow}}$ microseconds,
plus the fixed overhead of the AABB and hint-axis filters.
Chapter 4 uses the benchmark throughput numbers from
Table~\ref{tab:bench-results} to project feasible SA step rates for
each problem size $N$.